\section{n2. Explicit normal-crossing surface: signed primary-model obstacle has no greatest geometric C or D minorant.}
\textit{Supports: G331, n1.}

The proposed surface obstruction is valid for the ACTUAL G331 geometric C/D cone, over EVERY field k. It also has a stronger signed primary-model version and an exact canonical filtration. Put R=k[x,y,z]\_(x,y,z)/(xy), a=m=(x,y,z), X=V\_m, retaining ALL supports. This is a reduced excellent nonnormal surface. Its finite normalization is S=k[x,z]\_(x,z) times k[y,z]\_(y,z); R is the fibre product over k[z]\_(z). Pairs have identical x=0,y=0 residues; subtracting their common residue shows R is generated by z and the separated x- and y-terms. The component idempotents make S finite. Testing multiplication by those idempotents shows its conductor is exactly c=(x,y). The reduced conductor quotients are k[z]\_(z) and two copies of it. Their normalized boundary consists of two points over one point o: support(x,y), v\_o(z)=1. G331 identifies precisely these two source points, with singleton fibres everywhere else.

On branch1, y=0, set f1=-1/2. On branch2, x=0, set f2(w)=-1/2-1/(2w(y)), with 1/infinity=0. Normalization gives min(w(y),w(z))=1, so w(y)\textgreater{}=1. Extended evaluation w(y) is continuous, and reciprocal on [1,infinity] is continuous and tends to0. Hence f2 is continuous, in[-1,-1/2]. Both source values at EVERY conductor fibre are -1/2. The compact finite quotient topology gives a genuine continuous signed function f downstairs. No unsupported claim that f is psh is made.

Let v1 be the branch2 GENERIC monomial valuation with y=z=1; x is killed. For ANY proper primary ideal I, e=v\_o(I) is a positive finite integer. Choose an element attaining e and write it as a polynomial divided by a local unit. In its branch2 numerator the y=0 residue has a nonzero z\textasciicircum{}e coefficient. The generic monomial valuation has no cancellation between distinct monomials, so its total-degree minimum is\textless{}=e. The denominator has nonzero constant and value0 at both valuations. Thus v1(I)\textless{}=v\_o(I). Unit ideals give equality0. Negative combinations consequently satisfy q(o)\textless{}=q(v1). G331's proved uniform comparison applies to the ENTIRE geometric C, including real Cartier core evaluations and uniform closure, so EVERY g in C satisfies g(o)\textless{}=g(v1). Taking infima in an ACTUAL decreasing C-net proves the same inequality for EVERY g in D, including bottom. This is the substantive geometric bridge; it uses a supplied comparison theorem, not a new ideal-defined cone.

Every C/D minorant of f therefore has g(o)\textless{}=f(v1)=-1. Moreover EVERY normalized semivaluation on the regular second branch dominates its ordinary m-order: polynomials in m\textasciicircum{}e have value\textgreater{}=e and units have value0. By the same uniform and decreasing bridge, g(w)\textless{}=g(v1) for ALL branch2 w. Thus every such minorant is\textless{}=-1 throughout branch2, including arbitrary supports.

For n\textgreater{}=1 let I\_n=m\textasciicircum{}(2n)+x\textasciicircum{}2 m\textasciicircum{}n and h\_n=-v(I\_n)/(2n). This ideal is proper primary because it contains m\textasciicircum{}(2n) and is contained in m. Ideal sum/product identities hold on EVERY support. On branch1,
h\_n=max(-1,-1/2-v(x)/n),
whereas on branch2 and at o it is identically-1. On ACTUAL Bl\_(I\_n), I\_n O is invertible with a regular generator, defining effective Cartier E\_n; -E\_n/(2n) is negative effective REAL Cartier and relatively ample. Its degree is positive on EVERY contracted integral curve, before any limit is taken. Hence h\_n is an actual geometric-core competitor. It is below f everywhere. On every branch1 point with finite x-value, h\_n increases to-1/2; the sole branch1 point killing x is the conductor point. The constant-1 is also a competitor. The COMPLETE raw C/D supremum is consequently
p=-1/2 on branch1 minus o; p=-1 on branch2, including o.
The upper bound on branch1 is f1; the upper bound on branch2 was proved for every geometric competitor. This determines p on the entire boundary, not only on sampled monomial points.

Define branch1 w\_r(x)=r,w\_r(z)=1,w\_r(y)=infinity, r\textgreater{}=1, by generic monomial minima. For each polynomial, as r tends to infinity its minimum tends to the order of its x=0 residue, or infinity if that residue is0. All local unit denominators have value0. Thus w\_r converges pointwise on EVERY element of R to o. But p(w\_r)=-1/2 and p(o)=-1. The raw supremum is not USC or continuous. Every D function is USC, being an infimum of continuous functions. A greatest C or D minorant would equal the respective raw supremum (the two agree by n1), so neither exists. Equivalently such a greatest minorant must dominate all h\_n, take value-1/2 on the finite-x branch1 locus and\textless{}=-1 at o, contradicting USC. The USC regularization p* has p*(o)=-1/2, p*(v1)=-1; away from o the two component loci are open and p* equals p. It violates the geometric inequality. It remains\textless{}=f, but cannot be declared geometrically positive.

A finite signed primary-model obstacle has the SAME failure. Let J=(y)+m\textasciicircum{}2, an actual proper primary ideal, and F=(1/2)v(J)-(3/2)v(m). On X, branch1 gives F=-1/2; branch2 gives F=(1/2)min(w(y),2)-3/2 in[-1,-1/2]; F(o)=-1/2 and F(v1)=-1. This is a signed actual exceptional REAL Cartier evaluation on the directed common primary blowup of m and J. The SAME h\_n are below F, and the SAME universal inequalities prove exactly the same raw p and no greatest C/D minorant. The failure already occurs for a signed model obstacle.

For BOTH f and F the canonical ideal in EVERY degree m\textgreater{}=1 is explicitly
A\_m=(z\textasciicircum{}m)+x(x,z)\textasciicircum{}(ceil(m/2)-1)+y(y,z)\textasciicircum{}(m-1).
Indeed on branch1 its requirement is equivalent to ordinary order\textgreater{}=ceil(m/2), by the generic order valuation and domination by all normalized supports. On branch2 the requirement is equivalent to order\textgreater{}=m: necessity is v1, sufficiency follows from f2,F2\textgreater{}=-1. A pair with these orders belongs to R precisely when its common z-residue has order\textgreater{}=m. Separate the common residue and its x- and y-divisible remainders to obtain the displayed ideal, including arbitrary local-unit denominators. Thus on branch1, writing s=v(x),t=v(z), min(s,t)=1,
h\_m=-min(mt,s+ceil(m/2)-1)/m;
on branch2 and at o, h\_m=-1. This is an all-support ideal identity, not an experimentally inferred filtration.

All h\_m have actual negative effective ample primary presentations, and converge pointwise in EVERY degree to p. Nevertheless ||p-h\_m||\_infinity=1/2 for EVERY m: choose t=1 and finite s\textgreater{}=m-ceil(m/2)+1 to get h\_m=-1,p=-1/2; elsewhere the difference lies in[0,1/2]. The canonical algebra is not finitely generated, by n1's elementary forward implication from finite generation to a continuous rational-primary envelope. Here non-FG accompanies a proved failure; it is not the reason substituted for that failure.

The failure is stable on an open set of continuous obstacles. If ||f'-f||\textless{}epsilon\textless{}1/4, every h\_n-epsilon is a geometric C minorant of f'. A greatest C/D minorant would therefore have values\textgreater{}=-1/2-epsilon on the finite-x first-component locus. At o the universal two-point inequality forces value\textless{}=f'(v1)\textless{}=-1+epsilon. USC along w\_r would require the opposite lower bound -1/2-epsilon at o, impossible when2epsilon\textless{}1/2. This also shows that failure is not a single fragile equality case.

Homogeneity is explicit. On the unnormalized cone write lambda=v(m). The original branch2 obstacle is -lambda/2-lambda\textasciicircum{}2/(2v(y)), interpreted as- lambda/2 if v(y)=infinity; branch1 is-lambda/2. It is homogeneous of degree one. F already is a homogeneous signed ideal evaluation. The raw p is -lambda/2 on the finite-x first-component locus and-lambda otherwise. For a new primary gauge b divide the same homogeneous functions by v(b) and use G331's exact ray homeomorphism; bounded positive denominators preserve the discontinuity and the no-greatest conclusion. Canonical ideals remain identical under this transported obstacle. Completion also preserves the example.

This is a genuinely NONNORMAL envelope obstruction with a two-point infinite-conductor fibre. Each normal surface component individually has CE by n1. It contradicts neither the bijective-normalization theorem reconstructed in n1 nor reduced-curve min(f,0): curve conductors have finite value and no such identified boundary points. Literal classical/CLD radial profiles, products, energy and full finite-map operators remain separate. The new exact formulas require no new resolution, finite generation, F-finiteness or characteristic assumption. Artifacts n2.md and n2.tex; finite rational-weight diagnostics are in diagnostics.py and reports/diagnostics.json. Rerun python diagnostics.py; python deliver.py; python check\_delivery.py. The proofs above are universal; the finite checks are observations only.

\subsection*{Exact crossing formulas}
\[
 R=k[x,y,z]_{(x,y,z)}/(xy),\quad
 J=(y)+\mathfrak m^2,\quad F=\tfrac12v(J)-\tfrac32v(\mathfrak m).
\]
\[
 g(o)\le g(v_1)\quad(g\in C\cup D),\qquad
 h_n=-\frac{v(\mathfrak m^{2n}+x^2\mathfrak m^n)}{2n}.
\]
\[
 p=\begin{cases}-\tfrac12&\text{first component, }v(x)<\infty,\\
 -1&\text{otherwise},\end{cases}\qquad p\notin\operatorname{USC}(X).
\]
\[
 A_m=(z^m)+x(x,z)^{\lceil m/2\rceil-1}+y(y,z)^{m-1},\qquad
 \|p-h_m\|_\infty=\tfrac12\quad(m\ge1).
\]

